\chapter{PFN pro segmentaci obrazu}
\label{chap:segmentace}


V předchozí kapitole jsme PFN zkoumali jen na jednorozměrné regresi.
Tam jsme se měli o~co opřít, pro Gaussovský proces umíme spočítat přesný posterior.
V~této kapitole dáme PFN těžší a praktičtější úlohu, a to \emph{segmentaci obrazu}.
Ukážeme, že není to zcela nová úloha, ale přirozené pokračování celé práce.
Jde totiž o~stejnou amortizovanou Bayesovskou inferenci, jen o~jednu dimenzi výše a s~binárním cílem.
Současně tím získáme nové, nezávislé prostředí, ve kterém můžeme znovu ověřit předpovědi Naglerovy teorie ze sekce~\ref{sec:pfn-teorie}.

\section{Od jednorozměrné regrese k segmentaci obrazu}
\label{sec:seg-uvod}

V~této úvodní sekci popíšeme tuto novou úlohu a přeložíme do ní pojmy, se kterými jsme dosud pracovali.
Doposud PFN dostávalo body na přímce a předpovídalo spojitou hodnotu.
Segmentace je úloha jiného typu.
Na vstupu máme obrázek a chceme u~každého jeho bodu (pixelu) rozhodnout, zda patří do hledané oblasti nebo ne.
Tomuto rozhodování pro celý obrázek najednou se říká \emph{segmentace}, tedy rozdělení obrazu na oblast zájmu a pozadí.
Oblasti zájmu, tedy pixelům hledaného útvaru, budeme říkat \emph{popředí} (anglicky \emph{foreground}), všem ostatním pixelům \emph{pozadí} (\emph{background}).
Výsledkem je tzv. \emph{maska}, což je obrázek stejného rozměru jako vstup, ve kterém má pixel popředí hodnotu jedna a pixel pozadí hodnotu nula.
Masku si můžeme představit jako černobílý obrys obtažený přes původní obrázek.
\par

Hlavní rozdíl oproti regresi je ve výstupu.
V~regresi PFN předpovídalo rozdělení jedné spojité hodnoty pro jeden vstupní bod.
Zde předpovídá pro každý pixel pravděpodobnost, že leží uvnitř oblasti, tedy číslo mezi $0$ a $1$.
Model tedy nevrací rovnou tvrdou černobílou masku, ale \emph{pravděpodobnostní mapu}, kde hodnota blízká jedné znamená „téměř jistě uvnitř“ a hodnota blízká nule „téměř jistě mimo oblast zájmu“.
Právě tato mapa nás zajímá, protože u~ní můžeme stejně jako v~jedné dimenzi zkoumat přesnost predikce.
\par

Ukažme nyní, že segmentace není zcela nová úloha, ale přirozené pokračování toho, co už jsme viděli v~kapitole~\ref{chap:experimenty}.
V~celé předchozí kapitole jsme pracovali s~funkcí jedné proměnné vzorkovanou z~Gaussovského procesu.
Představme si nyní stejnou myšlenku, ale s~funkcí \textit{dvou} proměnných, tedy s~náhodným polem přes celou plochu obrázku.
Hladkost tohoto skrytého pole řídí stejná korelační délka $\ell$ jako dřív.
 Malé $\ell$ dává drobné a ostré útvary, velké $\ell$ velké hladké plochy.
Zatím je to jen dvourozměrná obdoba Gaussovského procesu z předchozí kapitoly.
Z~tohoto spojitého pole vytvoříme masku jednoduchým krokem.
Zvolíme práh a řekneme, že pixel patří do popředí právě tehdy, když hodnota pole v~něm práh překročí.
Tím se ze spojitého pole stane binární maska, tedy náš cíl.
Vstupní obrázek pak vznikne tak, že k~poli přidáme šum.
Obrázek je tedy rozmazaná verze skryté pravdy a model má z~něj odhadnout ostrou masku pod ní.
Obrázek~\ref{fig:seg-priklady} ukazuje dva takové páry obrázku a jeho masky, jeden pro velké $\ell$ (hladké skvrny) a jeden pro malé $\ell$ (drobné skvrny).
\par

\begin{figure}[!tb]
    \centering
    \includegraphics[width=0.72\linewidth]{figures/VariantC_pfn_seg/fig_seg_intro_examples.png}
    \caption{Ukázka úlohy segmentace. Vlevo vstupní obrázek (zašuměné skryté pole), vpravo odpovídající maska (cíl predikce), kterou získáme prahováním skrytého pole. Horní řádek odpovídá velké korelační délce $\ell$, tedy hladkým plochám, dolní řádek malé $\ell$, tedy drobným ostrým útvarům. Model dostává na vstupu levý obrázek a má předpovědět pravou masku.}
    \label{fig:seg-priklady}
\end{figure}

Zbývá vysvětlit, jak modelu řekneme, jaká pravidla pro danou sérii obrázků platí, tedy jaká je korelační délka, kde leží práh a jak silný je šum.
Stejně jako v~regresi mu to neřekneme přímo, ale ukážeme mu několik hotových příkladů.
Zde opět postupujeme podle standardního postupu.
Dáme mu sadu dvojic obrázku a už správně nakreslené masky a model si z~nich má odvodit, jaká pravidla v~této sérii platí.
Této sadě příkladů budeme říkat \emph{support set} (kontextová sada), je to přímá obdoba kontextových bodů $D_n$ z~jednorozměrného případu.
Model tedy funguje a zpracovává vstup podle stejného principu, jak jsme již popsali v~kapitole~\ref{chap:pfn}.
\par

Celkově tedy jde o~stejnou amortizovanou Bayesovskou inferenci, jen ve dvou dimenzích a s~binárním cílem.
Posteriorní prediktivní rozdělení (PPD), které bylo v~regresi Gaussovou hustotou, je zde v~každém pixelu Bernoulliho rozdělení, tedy pravděpodobnost, že pixel patří do oblasti.
Trénovací kritérium, kterým byla u~regrese záporná log-věrohodnost výstupní hlavy \texttt{BarDistribution}, se zde stává pixelovou binární křížovou entropií (BCE).
Všechny ostatní pojmy zůstávají beze změny, jen je čteme v~novém kontextu.
\par

Nyní můžeme říct, co přesně budeme v~této kapitole zkoumat.
Připomeňme, že Naglerovy výsledky ze sekce~\ref{sec:pfn-teorie} nezávisí na dimenzi vstupu.
Věta~\ref{theorem:Rozptyl_transformeru} o~mizení rozptylu i neredukovatelný bias z~podmínky lokality (věta~\ref{theorem:lokalizace}) jsou tvrzení o~\emph{attention mechanismu} obecně, a měla by tedy platit i pro segmentaci.
Právě tohoto faktu využijeme.
Nejprve vezmeme hotový, už natrénovaný \emph{in-context} segmenter UniverSeg~\cite{butoi2023}.
Není to PFN, ale funguje na stejném principu, predikuje z~kontextové sady přes globální \emph{attention}.
Ověříme, zda se u~něj objeví tzv. \emph{kolaps k~prioru}.
Jde o~jev, kdy model s~příliš malým kontextem přestane odpovídat na konkrétní vstup a vrátí se k~výchozí odpovědi, tedy k~průměrné masce, kterou zná z~tréninku.
Intuice je jednoduchá.
Když má model málo příkladů, nedokáže z~nich poznat pravidla úlohy, a nejbezpečnější volbou je pro něj průměr.
S podmínkou lokality tento jev souvisí následovně.
\emph{Softmax attention} nikdy nedá žádnému bodu kontextu váhu přesně nula.
Predikce proto vždy obsahuje složku tvořenou průměrem přes celý kontext, která na sousedech konkrétního bodu nezávisí.
Při dostatku dat se tato globální složka projevuje jen jako malý neredukovatelný bias.
Při malém nebo neinformativním kontextu z~predikce zbude právě jen zprůměrovaná odpověď známá z~tréninku.
Pokud se kolaps objeví i u~modelu, který jsme sami netrénovali, nejde o~artefakt našeho tréninku, ale o~obecnou vlastnost amortizované inference.
Poté natrénujeme vlastní PFN přímo na úloze popsané výše.
Zde využijeme hlavní výhodu syntetických dat, pravidla jejich generování známe, a proto umíme pro každý obrázek spočítat i přesný posterior.
Tomu budeme stejně jako v~jedné dimenzi říkat \emph{oracle}.
U~reálných segmentačních dat takovou pravdu nikdy nemáme.
Zde avšak proti ní můžeme přesně změřit, jak věrně ji PFN aproximuje a kde vzniká odchylka, tedy bias, rozptyl, kalibraci a chování mimo trénovací rozdělení.
Zkoumáme tedy, jak dobře PFN aproximuje ideální Bayesovskou inferenci a kde se od ní začíná odchylovat, jen nyní na náročnější úloze.

\section{Kolaps k prioru u předtrénovaného modelu}
\label{sec:seg-universeg}

V~tomto prvním experimentu nebudeme trénovat vlastní model, ale vezmeme hotový \emph{in-context} segmenter a otestujeme na něm předpověď Naglerovy podmínky lokality.
\par

\subsection{Uspořádání experimentu}
\label{subsec:seg-a-setup}

Jako model použijeme UniverSeg, což je předtrénovaný \emph{in-context} segmenter.
Model necháme zmražený, žádnou jeho váhu neupravujeme.
UniverSeg nabízí interní průměrování přes více kontextových sad.
Voláme ho vždy nad jedním \emph{support setem}, tedy bez tohoto interního průměrování.
Kdybychom průměrovali, zaměnili bychom vliv velikosti kontextu s~vlivem tohoto průměrování, a to nechceme.
\par

Kolaps budeme sledovat podél dvou os.
Hlavní a čistou osou je velikost \emph{support setu} $n_{\text{supp}}$, kterou měníme od jednoho příkladu po $64$.
Na ní stojí hlavní tvrzení experimentu, protože přímo řídí, kolik informace model dostane, a nezávisí na tom, o jaká data jde.
Druhou, doplňkovou osou je změna domény, tedy typu a zdroje obrázků.
Postupujeme přes tři domény, snímky bílých krvinek (WBC-JTSC~\cite{zheng2018}), jinou akvizici stejných buněk (WBC-CV~\cite{zheng2018}) a mozkové MRI (OASIS~\cite{marcus2007}).
Poznamenejme, že touto osou neměříme vzdálenost dat od tréninku, model všechny tři modality při tréninku viděl, ale jen robustnost napříč různými úlohami.
U krvinek je cílem segmentace maska celé buňky, u~mozku maska bílé hmoty.
\par

Abychom mohli měřit posun predikce k~prioru, potřebujeme prior nějak vyjádřit.
UniverSeg žádný explicitní prior nemá.
To obejdeme tímto způsobem.
Nahradíme ho operační prior maskou, tedy průměrnou maskou přes všechny příklady dané domény.
Ta má smysl jen pro prostorově zarovnaná data, kde průměr dává rozumný tvar.
U~krvinek je buňka vždy zhruba uprostřed, takže průměrná maska je kompaktní útvar uprostřed, jak ukazuje obrázek~\ref{fig:seg-a-prior}.
U~mozku je struktura složitější a zašuměnější, ale stále čitelná.
\par

\begin{figure}[!tb]
    \centering
    \includegraphics[width=\linewidth]{figures/VariantA_universeg/fig_A_05_prior_masks.png}
    \caption{Operační prior masky, tedy průměrné masky přes všechny příklady dané domény. Hodnota v~pixelu udává, jak často v~něm napříč doménou leží popředí. U~krvinek (WBC) je maska kompaktní útvar uprostřed, u~mozku (OASIS) je strukturovanější. Hodnota \emph{fg} je podíl popředí v~průměrné masce.}
    \label{fig:seg-a-prior}
\end{figure}

\subsection{Kolaps k prioru}
\label{subsec:seg-a-kolaps}

Chceme zjistit, zda se predikce s~ubývajícím kontextem vrací k~prioru, jak předpovídá věta~\ref{theorem:lokalizace}.
K~tomu zavedeme dvě vzdálenosti.
První je vzdálenost predikce od pravé masky, značíme ji $d_{\text{truth}}$.
Druhá je vzdálenost predikce od operační prior masky, značíme ji $d_{\text{prior}}$.
Obě počítáme jako průměrnou pixelovou kvadratickou odchylku.
Malé $d_{\text{truth}}$ tedy znamená, že model předpověděl blízko pravdě, malé $d_{\text{prior}}$, že se drží výchozí průměrné masky.
\par

\textbf{Diskuze:}

Kolaps se na obrázku~\ref{fig:seg-a-kolaps} pozná podle křížení obou křivek.
Při kontextu rovnajícím se jediné ukázce je predikce ve všech třech doménách blíž k prior masce než pravdě, model tedy hádá zhruba průměrnou masku.
S tím jak roste \emph{support set}, $d_{\text{truth}}$ prudce klesá a někde se protne s~$d_{\text{prior}}$, od té chvíle už je predikce blíž pravdě než prioru.
Právě tento příklon k~prioru při slabé podmiňující informaci Naglerova teorie předpovídá, a potvrzuje se tak i na modelu, který není úplně \emph{PFN-type}.
Poloha tohoto křížení se mírně liší podle domény, ale jev je ve všech třech stejný.
Kolaps tedy řídí velikost kontextu, nikoli konkrétní úloha.
Podstatné je také, že $d_{\text{prior}}$ neklesá k~nule, ale ustálí se na nenulové hodnotě.
Model totiž ani při velkém kontextu nepredikuje přesně průměrnou masku, jeho predikce je specifická pro daný obrázek.
\par

\begin{figure}[!tb]
    \centering
    \includegraphics[width=\linewidth]{figures/VariantA_universeg/fig_A_01_collapse_crossover.png}
    \caption{Kolaps k~prioru. Zelená křivka je vzdálenost predikce od pravé masky $d_{\text{truth}}$, fialová vzdálenost od operační prior masky $d_{\text{prior}}$, obojí jako funkce velikosti \emph{support setu}. Při jediném příkladu je predikce blíž prioru než pravdě (kolaps), s~rostoucím kontextem se křivky protnou a predikce přilne k~pravdě. Chybové úsečky jsou $\pm 1$ směrodatná chyba.}
    \label{fig:seg-a-kolaps}
\end{figure}

\subsection{Výkon a saturace kontextem}
\label{subsec:seg-a-dice}

Mimo jiné by nás také mohla zajímat i samotná kvalita masky.
Měříme ji \emph{Dice} koeficientem, tedy mírou překryvu predikované a pravé masky, kde jedna znamená dokonalý překryv a nula žádný.
\par

\textbf{Diskuze:}

Obrázek~\ref{fig:seg-a-dice} ukazuje, že \emph{Dice} s~rostoucím kontextem roste a poté \emph{saturuje}.
Saturací rozumíme stav, kdy křivka přestane růst a ustálí se na plató, další příklady v~kontextu už tedy kvalitu masky nezlepšují.
V~našem případě se přidávat další kontext přestane vyplácet zhruba od osmi až šestnácti příkladů.
Je to přímá dvourozměrná obdoba saturace přesnosti s~velikostí kontextu, kterou jsme viděli u~regrese v~kapitole~\ref{chap:experimenty}.
Plató se mírně liší podle domény, ale kvalitativní tvar křivky, tedy růst a následná saturace, zůstává stejný.
\par

\begin{figure}[!tb]
    \centering
    \includegraphics[width=0.72\linewidth]{figures/VariantA_universeg/fig_A_02_dice.png}
    \caption{\emph{Dice} koeficient jako funkce velikosti \emph{support setu} pro tři domény. Křivka roste a saturuje, výška plató se mírně liší podle domény. Chybové úsečky jsou $\pm 1$ směrodatná chyba.}
    \label{fig:seg-a-dice}
\end{figure}

\subsection{Kalibrace}
\label{subsec:seg-a-kalibrace}

Nakonec nás zajímá, zda pravděpodobnosti, které model vrací, odpovídají skutečnosti.
Řekne-li model u~nějakých pixelů hodnotu $0{,}8$, mělo by z~nich popředí obsahovat zhruba $80\,\%$, jinak jsou jeho pravděpodobnosti zavádějící.
Tuto shodu měříme dvěma metrikami a u~obou platí, že směr případné chyby nepředpokládáme, ale měříme.
První je \emph{reliability diagram}.
Vezmeme všechny pixely, u~kterých model hlásí zhruba stejnou pravděpodobnost popředí, například $\hat p \approx 0{,}7$, a spočítáme, jaký podíl z~nich v~pravé masce popředí skutečně je.
U~dokonale kalibrovaného modelu by vyšlo právě $70\,\%$.
Diagram tyto dvojice vynáší proti sobě, dokonalá kalibrace tedy leží na přímce.
Leží-li křivka nad přímkou, je skutečný podíl popředí vyšší, než model hlásí, a model tedy pravděpodobnost popředí podhodnocuje.
Odchylku od této přímky shrnuje jedno číslo, očekávaná kalibrační chyba (ECE).
Druhou metrikou je tzv. \textit{znaménkový rozdíl}, který hodnotí tvrdá rozhodnutí.
Pixel prohlásíme za popředí, pokud $\hat p > 0{,}5$, jinak za pozadí.
Každé takové rozhodnutí má svou jistotu, a tou je pravděpodobnost, kterou model zvolené třídě přiřadil, tedy $\max(\hat p,\, 1-\hat p)$.
Řekne-li model například $\hat p = 0{,}9$, rozhodne pro popředí s~jistotou $0{,}9$, řekne-li $\hat p = 0{,}2$, rozhodne pro pozadí s~jistotou $0{,}8$.
Znaménkový rozdíl je pak průměrná jistota minus podíl správných rozhodnutí.
Kladná hodnota znamená, že si model svými rozhodnutími věří víc, než jak často má pravdu, záporná naopak.
\par

\textbf{Diskuze:}

Měření na obrázcích~\ref{fig:seg-a-reliability} a~\ref{fig:seg-a-kalibrace} ukazuje, že směr kalibrace není univerzální, ale závisí na doméně.
U~krvinek můžeme říct, že je model podhodnocený.
Model je opatrný a vrací mírně jiné pravděpodobnosti, než by odpovídalo skutečnosti.
U~mozku je naopak přehnaně sebejistý, a to téměř nezávisle na velikosti kontextu.
Znaménkový rozdíl na obrázku~\ref{fig:seg-a-kalibrace} vpravo je přitom kladný ve všech doménách, nejsilnější je při jednom příkladu a s~rostoucím kontextem klesá k~nule.
To není ve sporu s~reliability diagramy, obě metriky totiž odpovídají na jinou otázku.
\emph{Reliability diagram} se ptá, zda sedí hlášená pravděpodobnost popředí, znaménkový rozdíl se ptá, zda sedí jistota zvolené třídy.
Model tak může u~krvinek pravděpodobnost popředí podhodnocovat, tedy mít křivku nad úhlopříčkou, a současně si svými tvrdými rozhodnutími věřit víc, než jak často má pravdu.
Když navíc oddělíme vnitřek objektu od úzkého pásma kolem hranice, ukáže se, že se kalibrační chyba soustředí právě na hranici.
Kalibrační chyba je tam řádově vyšší než uvnitř a s~kontextem klesá mnohem pomaleji, u~mozku dokonce zůstává prakticky konstantní.
Poznamenejme, že tento závěr je v~souladu se vzorem, který jsme viděli i v~jedné dimenzi, model dobře pozná, \textit{kde} je nejistý, ale absolutní míru nejistoty odhaduje hůře.
\par

\begin{figure}[!tb]
    \centering
    \includegraphics[width=\linewidth]{figures/VariantA_universeg/fig_A_03_reliability.png}
    \caption{\emph{Reliability} diagramy pro tři domény (řádky) a vybrané velikosti \emph{support setu} (sloupce). Křivka nad úhlopříčkou znamená podhodnocenou jistotu, pod úhlopříčkou přehnanou.}
    \label{fig:seg-a-reliability}
\end{figure}

\begin{figure}[!tb]
    \centering
    \includegraphics[width=\linewidth]{figures/VariantA_universeg/fig_A_04_calibration.png}
    \caption{Vlevo očekávaná kalibrační chyba (ECE) zvlášť pro celý obrázek a pro hraniční pásmo, vpravo znaménkový rozdíl mezi jistotou a správností (kladný znamená přehnanou jistotu). Kalibrační chyba je na hranici řádově vyšší než uvnitř a klesá s~kontextem pomaleji.}
    \label{fig:seg-a-kalibrace}
\end{figure}

\subsection{Shrnutí}
\label{subsec:seg-a-shrnuti}

Tento experiment ukázal, že hotový \emph{in-context} segmenter dědí přesně ty samé patologie, které Naglerova teorie předpovídá pro každý amortizovaný model s~globálním \emph{attention}.
S~ubývajícím kontextem se predikce vrací k~prioru a kalibrace je nestabilní.
Protože jsme model sami netrénovali, nemůže jít o~artefakt naší trénovací procedury, ale o~obecnou vlastnost amortizace.
Podstatné je, že jev je napříč všemi třemi úlohami kvalitativně stejný.
Kolaps tedy řídí velikost kontextu, nikoli konkrétní úloha.
\par

Současně se ukázala i mez tohoto testu.
Poznamenejme, že UniverSeg nemá pravý prior ani dostupný posterior, proto jsme se museli spolehnout na operační prior masku a nemohli jsme ji porovnat s~přesným posteriorem.
Přesně tuto mez odstraní následující experiment, ve kterém natrénujeme vlastní PFN na úloze, kde prior i přesný posterior známe, a stejné patologie změříme rigorózně proti \emph{oraclu}.




\section{Srovnání vlastního PFN s přesným posteriorem}
\label{sec:seg-pfn-oracle}

V~této sekci provedeme hlavní experiment kapitoly.
Natrénujeme vlastní PFN na úloze ze sekce~\ref{sec:seg-uvod}, ve které pro každý obrázek umíme spočítat přesný posterior.
Změříme, jak moc přesně PFN jej aproximuje a kde se od něj odchyluje.
\par

\subsection{Uspořádání experimentu}
\label{subsec:seg-c-setup}

Data generujeme přesně podle postupu z~úvodu kapitoly.
Vzorkujeme dvourozměrné pole z~Gaussovského procesu, prahováním z~něj vytvoříme masku a přidáním šumu vstupní obrázek.
Hyperparametry úlohy při tréninku nedržíme pevné.
Korelační délku, sílu šumu i podíl popředí (tedy jak velkou část obrázku zabírá oblast zájmu) losujeme pro každou dávku znovu, aby se model musel naučit celou rodinu úloh, ne jednu konkrétní.
Pole vzorkujeme na čtvercové mřížce $64 \times 64$ pixelů, obrázek má tedy stranu $L = 64$ pixelů.
Korelační délku $\ell$ přitom neudáváme v~pixelech, ale relativně, jako zlomek $\ell/L$ strany obrázku.
Takové vyjádření totiž přímo říká, jak velké jsou útvary vzhledem k~celému obrázku, a nezávisí na zvoleném rozlišení.
Například $\ell/L = 0{,}05$ znamená drobné útvary o~velikosti kolem pěti procent obrázku, $\ell/L = 0{,}40$ velké hladké plochy přes téměř polovinu obrázku.
Losujeme $\ell/L \in [0{,}05,\, 0{,}40]$ log-uniformně, směrodatnou odchylku šumu $\sigma \in [0{,}01,\, 0{,}30]$ a podíl popředí z intervalu $[0{,}15,\, 0{,}85]$.
Velikost \emph{support setu} pak volíme náhodně z množiny $\{1, 2, 4, 8, 16, 32, 64\}$.
Tyto intervaly tvoří trénovací rozsah, ke kterému budeme výsledky později vztahovat.
Jako architekturu záměrně volíme stejnou rodinu jako v~předchozí sekci, tedy UniverSeg, tentokrát ovšem s~náhodnou inicializací, bez předtrénovaných vah.
Model trénujeme od nuly na datech z~našeho prioru po dobu $250$ epoch, s~pixelovou Bernoulliho hlavou a s~binární křížovou entropií jako kritériem.
Právě tento trénink na explicitním prioru dělá z~modelu PFN.
Obě sekce se tak liší jen trénovací procedurou a daty, případné rozdíly v~chování proto nelze svést na architekturu.
\par

\textbf{Technická poznámka:}

Podstatná je jedna technická volba, díky které máme přesnou pravdu k~dispozici.
Výpočet přesného posterioru totiž vyžaduje inverzi kernelové matice, která má řádek a sloupec pro každý pixel, a její přímá inverze by proto byla příliš pomalá.
Proto pole vzorkujeme na \emph{toru}, tedy na mřížce se slepenými okraji, kde levý okraj sousedí s~pravým a horní s~dolním.
Na toru vypadá okolí každého pixelu stejně, takže kovariance dvou pixelů závisí jen na jejich vzájemném posunu.
Takovou matici diagonalizuje Fourierova transformace a inverze se rozpadne na obyčejné dělení frekvenci po frekvenci, přesný posterior proto spočítáme rychle a bez aproximace.
Výsledek je náš \emph{oracle} z~úvodu kapitoly, tedy nejlepší možná pravděpodobnostní mapa, jakou lze z~daného obrázku vůbec získat.
Torus je jen pomůcka pro tento výpočet a na interpretaci výsledků nemá podstatný vliv.
\par

Nakonec zaveďme značení režimů, podle kterého budeme výsledky číst.
Náročnost úlohy řídí korelační délka, čím je kratší, tím členitější je pole a tím těžší jeho predikce.
Uvnitř trénovacího rozsahu proto rozlišujeme úlohy snadné (\emph{Easy}, $\ell/L = 0{,}40$), střední (\emph{Medium}, $\ell/L = 0{,}15$) a náročné (\emph{Hard}, $\ell/L = 0{,}05$).
Snadný a náročný režim tak leží na okrajích trénovacího rozsahu.
K~nim přidáme dva režimy mimo trénovací rozsah, ve kterých je korelační délka kratší (\emph{OOD-short}, $\ell/L = 0{,}03$, členitější pole) nebo delší (\emph{OOD-long}, $\ell/L = 0{,}60$, hladší pole), než jakou model při tréninku viděl.
Zkratka OOD znamená \emph{out-of-distribution}, tedy data mimo trénovací rozdělení.
Není-li uvedeno jinak, každý vykreslený bod průměrujeme přes $40$ až $48$ nezávisle vygenerovaných úloh.
\par

\subsection{Shoda predikce s oraclem}
\label{subsec:seg-c-fidelita}

\begin{figure}[p]
    \centering
    \includegraphics[width=\linewidth]{figures/VariantC_pfn_seg/fig_C_01_qualitative.png}
    \caption{Ukázka predikcí pro snadnou (nahoře) a náročnou (dole) úlohu. Zleva vstupní obrázek, pravá maska, \emph{oracle} (pravý posterior), predikce PFN a jejich absolutní rozdíl. Mapy PFN a \emph{oraclu} téměř splývají, rozdíl je soustředěn na hranici mezi popředím a pozadím, tedy na obrysu segmentované oblasti.}
    \label{fig:seg-c-kvalita}

    \vspace{1.5em}

    \includegraphics[width=0.9\linewidth]{figures/VariantC_pfn_seg/fig_C_02_fidelity.png}
    \caption{Průměrná vzdálenost predikce PFN od \emph{oraclu} jako funkce velikosti kontextu pro pět režimů náročnosti. Uvnitř trénovacího rozsahu (Hard, Medium, Easy) je vzdálenost malá a na velikosti kontextu téměř nezávisí. Mimo rozsah (OOD-short, OOD-long) je větší, členitější pole (OOD-short) se \emph{oraclu} nepřiblíží ani s~velkým kontextem.}
    \label{fig:seg-c-fidelita}
\end{figure}

Nejprve se ptáme, jak blízko je predikce PFN k pravému posterioru.
Vzdálenost měříme absolutní odchylkou obou pravděpodobnostních map zprůměrovanou přes všechny pixely.
Na obrázku~\ref{fig:seg-c-kvalita} jsou vedle sebe vstupní obrázek, pravá maska, \emph{oracle} a predikce PFN pro snadnou a náročnou úlohu.
Mapy PFN a \emph{oraclu} jsou na pohled téměř k~nerozeznání.
Viditelný rozdíl zůstává jen na tenkých hranicích mezi popředím a pozadím.
\par

\textbf{Diskuze:}

Rozdíl mezi PFN a \emph{oraclem} znázorňuje graf na obrázku~\ref{fig:seg-c-fidelita}, který vynáší jejich průměrnou vzdálenost jako funkci velikosti kontextu pro všech pět režimů.
Plyne z~něj poznatek, který se v~této sekci bude opakovat.
Velikost kontextu téměř nehraje roli, predikce je dostatečně dobrá i s~jediným \emph{support párem}, tedy s~jedinou dvojicí obrázku a masky v~kontextu.
Vysvětlení je přímočaré a plyne přímo z~toho, jak naše úloha vzniká.
Připomeňme, že vstupní obrázek jsme vytvořili tak, že jsme ke skrytému poli, tedy k~té hladké funkci, ze které se prahováním dělá maska, přičetli šum.
Hodnota každého pixelu je proto hodnota tohoto pole v~daném bodě plus šum, tedy jeho přímé, jen zašuměné měření.
Obrázek jako celek tak zachycuje pole ve všech bodech najednou.
To je zásadní rozdíl oproti jednorozměrné regresi, kde model o~funkci věděl jen z~několika kontextových bodů.
Zde si model dokáže pole z~velké části zrekonstruovat už ze samotného vstupního obrázku, bez pomoci kontextu.
Ze \emph{support setu} se dozvídá jen to, co ze samotného obrázku poznat nejde, tedy kde leží práh mezi popředím a pozadím a jak silný je šum.
Co naopak rozhoduje, je náročnost úlohy.
Uvnitř trénovacího rozsahu je vzdálenost od \emph{oraclu} malá, mimo rozsah roste.
Nejhůře dopadají pole s~kratší korelační délkou, než jakou model viděl při tréninku.
U~nich chyba zůstává velká i s~velkým kontextem.
Naopak u~hladších polí se model přiblíží \emph{oraclu}, jakmile dostane malý kontext.
Stejnou asymetrii extrapolace jsme u~regrese viděli v~části~\ref{subsec:exp7-ood}, i tam byly krátké korelační délky pro model těžší než dlouhé.
\par

\subsection{Kalibrace}
\label{subsec:seg-c-kalibrace}

To, že střední hodnota vychází smysluplně, ještě nemusí nutně znamenat, že PFN správně reprezentuje i nejistotu.
Proto se ptáme, zda se pravděpodobnosti, které PFN vrací, shodují s~pravděpodobnostmi \emph{oraclu}, především na nejednoznačných hranicích.
\par

\textbf{Diskuze:}

Obrázek~\ref{fig:seg-c-kalibrace} vlevo ukazuje, že rozdělení predikovaných pravděpodobností PFN a \emph{oraclu} se téměř překrývají.
Obě rozdělení mají výrazné vrcholy u~nuly a u~jedné, tedy tam, kde je pixel jistě mimo nebo jistě uvnitř.
Mezi vrcholy leží nízké údolí.
Prostřední hodnoty pravděpodobnosti má jen malý zlomek pixelů.
To přesně odpovídá geometrii úlohy.
Nejisté jsou jen pixely v~úzkém pásmu podél hranic mezi popředím a pozadím, a těch je oproti vnitřkům oblastí málo.
Že jde právě o~hraniční pixely, přitom nevyčteme z~histogramu, ale z~konstrukce úlohy.
Prostřední pravděpodobnost dostane pixel právě tehdy, když je odhad pole blízko prahu, a body, kde pole protíná práh, tvoří z~definice masky právě její hranici.
Stejné pásmo prostředních hodnot podél hranic je ostatně přímo vidět na mapách na obrázku~\ref{fig:seg-c-kvalita}.
Podstatné je, že PFN tuto nejistotu na hranicích zachycuje stejně jako \emph{oracle}, nesnaží se ji uměle potlačit.
Zbývá jen mírná odchylka, na hranicích je totiž PFN o~něco sebejistější, než by odpovídalo \emph{oraclu}, jak ukazuje pravá část obrázku.
Tato přehnaná ostrost je jediná systematická vada kalibrace, kterou jsme naměřili, ale je malá.
Připomeňme, že jde o~stejný vzor jako v~jedné dimenzi, model dobře pozná, kde je nejistota, a jen mírně chybuje v~její velikosti.
\par

\begin{figure}[!htb]
    \centering
    \includegraphics[width=\linewidth]{figures/VariantC_pfn_seg/fig_C_05_calibration.png}
    \caption{Vlevo rozdělení predikovaných pravděpodobností PFN a \emph{oraclu} přes všechny pixely, obě křivky se téměř překrývají. Vpravo srovnání průměrné predikce PFN s~\emph{oraclem} po intervalech, mírné vybočení od úhlopříčky odpovídá přehnané ostrosti na hranicích.}
    \label{fig:seg-c-kalibrace}
\end{figure}

\subsection{Rozklad chyby na bias a rozptyl}
\label{subsec:seg-c-bias-var}

Nyní výsledky napojíme přímo na teorii ze sekce~\ref{sec:pfn-teorie}.
Chybu predikce rozložíme na dvě části.
První je rozptyl, tedy kolísání predikce podle toho, které konkrétní \emph{support páry} model dostal.
Druhou je bias, tedy systematická odchylka, která zůstane i po zprůměrování přes všechny volby kontextu.
Měříme je tak, že pro pevný vstupní obrázek opakovaně losujeme \emph{support set}, celkem $24$ tahů pro každou velikost kontextu, a sledujeme, jak se predikce mění.
Pro obě části máme z~Naglerovy teorie předpověď, kterou jsme v~jedné dimenzi potvrdili, a nyní ji ověříme znovu ve dvou.
\par

\textbf{Diskuze:}

Obrázek~\ref{fig:seg-c-biasvar} ukazuje obě části zvlášť.
Rozptyl s~rostoucím kontextem prudce klesá.
Připomeňme, že věta~\ref{theorem:Rozptyl_transformeru} omezuje odchylku predikce od jejího průměru řádem $n^{-1/2}$.
Rozptyl je střední hodnota kvadrátu této odchylky, a kvadrát $n^{-1/2}$ je $n^{-1}$.
Rozptyl by tedy měl mizet alespoň jako $n^{-1}$.
Pokles odečtený z~grafu tomuto řádu zhruba odpovídá.
U~některých úloh se zdá být i o~něco rychlejší, což horní mez věty neporušuje.
Přidávání dalších \emph{support párů} tedy skutečně ustaluje predikci.
Bias se naopak drží na malé kladné hodnotě a s~kontextem nemizí.
To je přesně neredukovatelný bias, který plyne z~podmínky lokality (věta~\ref{theorem:lokalizace}).
Jak jsme již popisovali v kapitole \ref{chap:experimenty}, globální \emph{attention} přiděluje váhu i vzdáleným bodům, takže drobná systematická odchylka nikdy úplně nezmizí.
\par

Z~jediného tréninku bychom ale nemohli tvrdit, že je bias vlastností architektury.
Mohl by to být jen zbytek konkrétního běhu.
Proto jsme natrénovali tři modely s~různou inicializací a různými trénovacími daty a bias změřili na stejné sadě úloh.
Obrázek~\ref{fig:seg-c-seedy} ukazuje, že se hodnota biasu napříč běhy shoduje na jednotky procent.
Tímto usuzujeme, že neredukovatelný bias je strukturální vlastnost architektury, a ne náhoda jednoho modelu.
Máme konkrétní praktické potvrzení jevu, který Naglerova teorie přesně předpovídá.
\par

\begin{figure}[!htb]
    \centering
    \includegraphics[width=\linewidth]{figures/VariantC_pfn_seg/fig_C_04_bias_variance.png}
    \caption{Rozklad chyby na bias a rozptyl jako funkce velikosti kontextu pro snadnou a náročnou úlohu (logaritmické osy). Vlevo plné čáry jsou bias, který zůstává na plató, přerušované rozptyl, který klesá. Vpravo rozptyl proti teoretickému řádu $n^{-1}$ plynoucímu z~věty~\ref{theorem:Rozptyl_transformeru}, empirický pokles tomuto řádu zhruba odpovídá.}
    \label{fig:seg-c-biasvar}
\end{figure}

\begin{figure}[!htb]
    \centering
    \includegraphics[width=0.72\linewidth]{figures/VariantC_pfn_seg/fig_C_07_seed_bias.png}
    \caption{Hodnota biasu na plató pro tři nezávislé tréninky (různá inicializace i data) na stejné sadě úloh. Body se shodují na jednotky procent, chybová úsečka je směrodatná odchylka přes běhy. Shoda ukazuje, že bias je strukturální vlastnost architektury.}
    \label{fig:seg-c-seedy}
\end{figure}

\FloatBarrier

\subsection{Vzdálenost od Bayesova optima}
\label{subsec:seg-c-bayes}

Dosud jsme PFN porovnávali s~\emph{oraclem} přímo.
Nyní se budeme dívat z jiné perspektivy, a to jak daleko je PFN od nejlepšího možného prediktoru.
I \emph{oracle} totiž dělá nenulovou chybu.
Obrázek je zašuměný, takže u~pixelů na hranici nejde s~jistotou určit, zda patří do popředí.
Existuje tedy nejmenší možná chyba, pod kterou se nedostane žádný prediktor, budeme jí říkat \emph{Bayes floor}.
Chybu měříme binární křížovou entropií (BCE) vůči pravé masce, pro PFN i pro \emph{oracle} stejně.
Rozdíl mezi chybou PFN a touto mezí je cena, kterou platíme za to, že PFN úlohu jen odhaduje z~kontextu.
\par

Popišme nejprve, co obrázek~\ref{fig:seg-c-bayes} zobrazuje a jak ho číst.
Na vodorovné ose obou grafů je pět režimů náročnosti, které jsme zavedli v~části~\ref{subsec:seg-c-setup}.
Zleva stojí OOD-short, tam je korelační délka kratší, než jakou model viděl při tréninku, a pole je proto členitější.
Následují tři režimy uvnitř trénovacího rozsahu, Hard, Medium a Easy, seřazené od nejkratší korelační délky po nejdelší.
Vpravo stojí OOD-long, tam je korelační délka naopak delší a pole hladší.
Levý graf ukazuje pro každý režim dva sloupce.
Modrý sloupec je chyba samotného \emph{oraclu}, tedy \emph{Bayes floor}, nejmenší chyba, jaké lze v~daném režimu vůbec dosáhnout.
Červený sloupec je chyba PFN.
Čím menší je rozdíl mezi nimi, tím blíže je PFN nejlepšímu možnému prediktoru.
Pravý graf vynáší tento rozdíl, tedy přebytečnou chybu (angl. \emph{excess risk}), přímo.
\par

\textbf{Diskuze:}

Nyní k~výsledkům.
Uvnitř trénovacího rozsahu (sloupce Hard, Medium a Easy) je přebytečná chyba malá.
U~režimů Hard a Medium sloupce PFN a \emph{Bayes floor} téměř splývají, u~režimu Easy je rozdíl viditelný, ale je stále malý.
Model se tedy nenaučil jen dobrý prediktor, ale téměř nejlepší možný.
Poznamenejme, že tento výsledek zpětně vysvětluje i jedno pozorování z~tréninku, které jsme dosud nekomentovali.
Trénovací chyba se totiž ke konci tréninku ustálila na malé, ale nenulové hodnotě.
Nebylo to tím, že by model uvázl, ale tím, že dosedl těsně nad \emph{Bayes floor}, pod který se žádný prediktor nedostane.
\par

Za povšimnutí stojí, že uvnitř rozsahu je přebytečná chyba největší právě u~hladkých polí (Easy), jeho sloupec je mezi režimy uvnitř rozsahu nejvyšší.
Graf pouze ukazuje tento fakt, příčinu z~něj nevyčteme.
Vysvětlení ale nabízí spojení dvou dřívějších pozorování.
Z~části~\ref{subsec:seg-c-kalibrace} víme, že PFN je na nejistých pixelech na hranicích o~něco ostřejší než \emph{oracle}.
A~hladká pole mají nejisté hranice nejširší, takže pixelů, na kterých se tato vada projeví, je v~nich nejvíc.
Součet mnoha malých chyb přes širokou hranici pak dává největší přebytečnou chybu.
Mimo trénovací rozsah (krajní sloupce OOD-short a OOD-long) se oba směry chovají různě.
Hladší pole (OOD-long) zůstávají srovnatelná s~režimem Easy, kdežto u~členitějších polí (OOD-short) přebytečná chyba vzroste na několikanásobek, tedy přesně tam, kde amortizace selhává.
\par

\begin{figure}[!htb]
    \centering
    \includegraphics[width=\linewidth]{figures/VariantC_pfn_seg/fig_C_06_bayes_floor.png}
    \caption{Vlevo chyba PFN a \emph{Bayes floor} (chyba samotného \emph{oraclu}) pro jednotlivé režimy, měřená binární křížovou entropií (BCE) na pixel. Vpravo jejich rozdíl $\mathrm{BCE}_{\mathrm{PFN}}-\mathrm{BCE}_{\mathrm{oracle}}$, tedy přebytečná chyba PFN nad optimem. Uvnitř trénovacího rozsahu je malá, prudce roste jen u~členitějších polí mimo rozsah (OOD-short).}
    \label{fig:seg-c-bayes}
\end{figure}

\FloatBarrier

\subsection{Kolaps k prioru nenastává}
\label{subsec:seg-c-kolaps}

V~sekci~\ref{sec:seg-universeg} jsme viděli kolaps k~prioru, tedy že se predikce při malém kontextu vrací k~průměrné masce.
Zde stejné měření provedeme na vlastním PFN, a to s~pravým priorem místo operační prior masky.
Prior masku totiž tentokrát známe přesně.
Před pozorováním vstupního obrázku patří každý pixel do popředí se stejnou pravděpodobností, rovnou podílu popředí.
Prior maska je tedy jednobarevná plocha bez jakéhokoli tvaru.
Měření provedeme na jednom zástupci uvnitř trénovacího rozsahu, konkrétně na náročné úloze (Hard).
Sledujeme vzdálenost predikce od \emph{oraclu} a od této prior masky jako funkci velikosti kontextu.
\par

I v tomto případě není zcela zřejmé, co zobrazuje graf na obrázku \ref{fig:seg-c-kolaps} a jak ho správně číst.
Proto se pokusíme opět čtenáři přiblížit co je vyneseno na tomto grafu.
Zelená křivka je vzdálenost predikce od \emph{oraclu}, fialová křivka je vzdálenost od uniformní prior masky (náš prior v tomto případě), obě jako funkce velikosti kontextu.
\par

\textbf{Diskuze:}

Podle popisu grafu výše můžeme udělat následující závěr.
Kdyby kolaps k prioru nastával i zde, viděli bychom fialovou křivku klesat k~nule při malém kontextu a zelenou naopak stoupat.
Predikce by se totiž podobala prior masce, a ne skutečnosti.
Avšak nic takového se zde neděje.
Zelená křivka zůstává nízko (blízko nuly) a nemění se s~velikostí kontextu, dokonce se drží u~nuly už i při jediném \emph{support páru}.
Predikce je tedy po celou dobu mnohem blíž \emph{oraclu} než prior masce.
Kolaps zde tedy nenastává.
Důvod je stejný jako v~části~\ref{subsec:seg-c-fidelita}, hustý vstupní obrázek prozradí úlohu hned u jednoho příkladu, takže model nemá důvod, proč se vracet k~průměru.
\par

Zdůrazněme, že tento negativní výsledek Naglerově teorii neodporuje.
Vyžaduje ale přesnější čtení toho, co teorie předpovídá.
Jelikož je PFN trénováno aproximovat PPD, přiblížení k~prioru je třeba čekat tam, kde je samotný posterior blízko prioru, tedy tam, kde je podmiňující informace o~úloze slabá.
V~naší úloze však nese vstupní obrázek zašuměné pozorování skrytého pole v~každém pixelu, tedy přímo tu veličinu, kterou prahujeme.
Úlohu proto z~velké části řeší už samotný vstupní obrázek a \emph{support set} dodává jen práh a sílu šumu (viz část~\ref{subsec:seg-c-fidelita}).
Posterior je tak daleko od prioru i při jediném \emph{support páru} a žádný kolaps proto teorie nepředpovídá, což naše měření potvrzuje.
Právě v~tom se tato úloha liší od experimentu se zmrazeným modelem v~sekci~\ref{sec:seg-universeg}.
Tam je vstupní obrázek sice rovněž hustý, ale sám o~sobě neurčuje, co se má segmentovat.
Stejnou buňku lze segmentovat po jádru, cytoplazmě i celé buňce a cílovou strukturu definuje až \emph{support set}.
Je-li \emph{support set} řídký, je řídký i jediný nosič informace o~úloze.
Posterior přes možné úlohy zůstává široký a nejbezpečnější volbou je průměrná maska.
Kdyby naše PFN přesto kolabovalo, nebylo by to potvrzení teorie, ale selhání aproximace posterioru.
Stopa příklonu k~prioru nicméně zůstává i zde.
Je jí malé bias plató z~části~\ref{subsec:seg-c-bias-var}, které s~kontextem nemizí a jehož neodstranitelnost vysvětluje právě podmínka lokality (věta~\ref{theorem:lokalizace}).
Podobný mechanismus, který u~řídkého a málo informativního kontextu způsobí kolaps, se u~informativního vstupu projeví jen touto drobnou systematickou odchylkou.
Kolaps k~prioru tedy není univerzální vlastnost amortizace.
Nastává tam, kde je jediným zdrojem informace o~úloze řídký kontext.
\par

\begin{figure}[!htb]
    \centering
    \includegraphics[width=0.72\linewidth]{figures/VariantC_pfn_seg/fig_C_03_collapse.png}
    \caption{Vzdálenost predikce PFN od \emph{oraclu} (zelená, $p_{\text{oracle}}$) a od uniformní prior masky (fialová, $p_{\text{prior}}$, tedy konstantní podíl popředí) jako funkce velikosti kontextu pro náročnou úlohu. Predikce je po celou dobu blízko \emph{oraclu} a daleko od prioru, ke kolapsu nedochází ani při jediném \emph{support páru}.}
    \label{fig:seg-c-kolaps}
\end{figure}

\FloatBarrier

\subsection{Shrnutí}
\label{subsec:seg-c-shrnuti}

Tento experiment ukázal, že jádro myšlenky PFN funguje i ve dvou dimenzích.
Skutečné PFN natrénované na explicitním prioru dostatečně dobře aproximuje pravý Bayesovský posterior, a to ve střední hodnotě i v~nejistotě.
Uvnitř trénovacího rozsahu je dokonce téměř Bayesovsky optimální.
Od nejlepšího možného prediktoru jej dělí jen malá chyba.
Rozklad chyby přitom odpovídá Naglerově teorii, rozptyl s~kontextem mizí alespoň jako $n^{-1}$ a bias zůstává na malé, ale strukturální hodnotě, a to bylo potvrzené shodou přes tři nezávislé tréninky.
\par

Současně se ukázalo, kde má metoda meze.
Amortizační chyba neroste s~ubývajícím kontextem, ale s~náročností úlohy a s~odchýlením hyperparametrů mimo trénovací rozsah.
Hustý vstupní obrázek totiž úlohu prozradí i s~malým počtem příkladů.
Oba experimenty této kapitoly tak vedou ke stejnému závěru.
Kolaps k~prioru u~předtrénovaného modelu ani selhání mimo trénovací rozsah nejsou vadou architektury, v~čistém režimu je totiž model prokazatelně téměř optimální.
Jsou to patologie podmíněné režimem, které nastupují až při řídkém kontextu nebo mimo naučenou rodinu úloh.
